% 本文件由 LaTeX 排版 Agent 根据有效 translation.json 自动生成。
% author=Gregory of Nyssa; work_id=2905; title=论并非三位天主

\chapter{论\Quote{并非三位天主}}

% structure: p0001 source=h1 target=omitted reason=page-title-already-rendered-by-parent

致\Person{阿布拉比乌}。\par

你既然在内里之人中充满全能的力量，理应有责任承担起对抗真理之敌的战斗，不应退缩于任务，使我们这些父亲因儿子们的崇高辛劳而喜乐；因为这是自然律的推动：但当你调整你的队列，将那些反对真理者投掷的标枪的冲击送向我们，并要求我们这些老年人用信德的盾牌扑灭他们的\Quote{灼热的火炭}和那被称为虚假的知识所磨利的箭矢时，我们接受你的命令，并使自身成为顺从的榜样，以便你本人能在类似的命令上给予我们公正的回报，基督的勇敢士兵\Person{阿布拉比乌}，若我们曾召唤你参与这样的竞赛。\par

事实上，你向我们提出的问题并非小事，也不是说若处理不当只会带来微小危害。因为问题的力量起初迫使我们接受两种错误意见之一，要么说\Quote{有三位天主}——这是不合法的，要么不承认圣子和圣神的天主性——这是不敬且荒谬的。\par

你陈述的论证大致如下：\Person{伯多禄}、\Person{雅各伯}和\Person{若望}，同在一个人性中，被称为三个人；对于在自然中合一者，若多于一人，用源自其本性的名称的复数来称呼他们，并无荒谬之处。那么，在上述情形中，既然习俗容许这样，无人禁止我们将两人称为两人，或两人以上称为三人，那么在我们陈述信仰奥秘时，虽然承认三个位格，也承认他们之间在本质上没有区别，但在某种意义上我们却与自己的宣认相矛盾：我们说父、子、圣神的天主性是单一的，却又禁止人们说\Quote{有三位天主}？这个问题，如我所说，非常棘手：然而，若我们能找到任何有助于支撑我们心灵的不确定性的东西，使之不再在这可怕的困境中摇摆不定，那将是好事；另一方面，即便我们的理性不足以解答此问题，我们也必须永远坚定不动摇地持守我们从教父们传承下来的传统，并从主那里寻求我们信仰的辩护理由：若这理由被任何蒙受恩宠者找到，我们必须感谢那赐予恩宠者；若找不到，我们仍将毫不动摇地持守那些已确定的论点。\par

那么，为何当我们在同一本性中逐一数算那向我们呈现的个体时，通常用复数称呼他们，说\Quote{如此多的人}，而不将他们全部称为一；而在神圣本性上，我们的教义定义却拒绝众神的复数，既枚举位格，同时又不承认复数含义？也许有人可能认为，若他（直率地对诚实的人）说：定义拒绝以任何数目计算神，是为了避免与异教徒的多神论有任何相似之处，以免我们若也像他们通常所做的那样不以单数而是以复数数算神性，可能被认为在某些教义上有共通之处。这个回答，我说，若对更纯朴的人作出，可能有些分量；但对于那些要求确立他们提出的两个选项之一的人（要么我们不承认三个位格中的神性，要么若承认，就应将共享同一神性的那些称为三个），这个回答不足以提供任何解决难题的方法。因此，我们必须更详尽地作出答复，尽我所能探究真理；因为这个问题非同寻常。\par

那么，我们首先指出，按照其共同本性的名称，用复数来称呼那些在本性上不可分割的人，说他们是\Quote{许多人}，是一种习惯性的语言滥用；而说他们是\Quote{许多人的本性}也几乎是一样的。我们可从以下例子看出此事的真理。当我们向任何人说话时，我们不用他本性的名称来称呼他，以免因名称的共通性引起混淆（假如每位听到的人都以为自己是那位被称呼的人，因为称呼不是用专属的名称，而是用他们本性的共通名称）；相反，我们藉着属于他本人的名称——即标识特定主体的名称——将他与群众分开。因此，有许多分享这本性的人——例如许多门徒、宗徒或殉道者——但在他们众人中，那个\OriginalTerm{人}{man}却是一；因为，如前所述，\Quote{人}这个术语并不属于某个个体的本性，而是属于共同的东西。因为\Person{路加}是\OriginalTerm{人}{man}，或\Person{斯德望}是\OriginalTerm{人}{man}；但若某人是人，却并不意味着他因此就是路加或斯德望。然而，位格的概念容许因各自考虑的独特属性而分开，当这些属性结合时，便藉数目呈现给我们；但他们的本性却是一，在自身内合一，是绝对不可分割的单位，不能借加增而扩大或因减除而缩小，在其本质上是且始终是一，即使表现在多重性中也是不可分离的、连续、完整，不与分享它的个体分裂。正如我们在每一种情形下都用单数谈论一个民族、或一群民众、或一支军队、或一个集会，而每一者都被视为复数，同样，按照更精确的表达，\Quote{人}应被称为一，即使那些以同一本性显现于我们面前的人构成了复数。因此，纠正我们错误的习惯，不再将本性的名称扩展到复数，远胜于因受制于习惯，而将上述情形中的错误转移到我们关于天主的陈述上。但既然纠正习惯不可行（因为谁能说服任何人不要将那些在同一本性中显现的人说成\Quote{许多人}呢？——事实上，在任何情形下，习惯都是难以改变的），我们在涉及低级本性的情况下不违背流行习惯，并不算太错，因为名称的误用不会带来危害；但在涉及天主本性的陈述中，术语的多种用法就不再如此安全了：因为在这些问题上，微不足道之事也不再是小事了。因此我们必须明认唯一天主，按照圣经的见证：\BibleQuote{以色列！你要听，上主你的天主是唯一的上主}，即使\OriginalTerm{天主性}{Godhead}的名称贯穿整个\OriginalTerm{天主圣三}{Holy Trinity}。我这么说，是按照我们在人性情形中所作的论述，从中我们得知，用复数的标记来扩展本性的名称是不恰当的。然而，我们必须更仔细地考察\Quote{\OriginalTerm{天主性}{Godhead}}这个名称，以便藉此词所包含的意义，获得一些帮助来澄清我们面前的问题。\par

大多数人认为\Quote{神性}一词在性质上被特别使用；正如天、太阳或宇宙的其他组成部分都由专有名称来指称对象一样，他们说，对于至高无上的神圣性质，\Quote{神性}一词恰当地适用于它向我们呈现的东西，作为一种特殊的名称。\newline{}我们则遵循圣经的提示，知道那性质是无可名状、不可言说的，并且我们说，无论是人类习俗所创造，还是圣经传授给我们的每一个术语，确实解释了我们关于神圣性质的观念，但并不包含那性质本身的意义。\newline{}这一点不难证明。因为用于受造物的所有其他术语，即使不分析其起源，也可发现它们是偶然地应用于对象，因为我们满足于用任何贴切的词语来指称事物，以避免对所意指之物的认识产生混淆。\newline{}但所有用来引导我们认识天主的术语，每个都包含其自身的含义，你找不到任何一个专门用于天主的词语是没有明确意义的。因此，很明显，通过我们使用的任何术语，并不表示神圣性质本身，而是揭示其某种属性。\newline{}例如，我们可能说天主是不朽的、大能的，或我们惯常对祂所说的任何其他话。但在这些术语中，每个都发现一种特殊的含义，适合理解或断言为神圣性质，但并不表达那性质本质是什么。因为主体，无论它是什么，是不朽的：但对不朽性的概念是这样的——那存在的东西不会分解为腐朽；所以，当我们说祂是不朽的时，我们宣告祂的性质不遭受什么，但并不表达那不受腐朽的是什么。\newline{}同样，如果我们说祂是生命的赐予者，虽然通过这个称呼我们表明祂给予什么，但我们并不通过这个词宣告那给予的是什么。\newline{}通过同样的推理，我们发现，所有其他从表达神圣属性的名称所涉及的含义中得出的结论，要么禁止我们想象那些不该想象的神圣性质的东西，要么教导我们那些应该想象的东西，但并不包含对性质本身的解释。\newline{}因此，既然我们感知到我们之上能力的各种运作，我们就从我们所知道的几种运作中塑造我们的称呼，并且我们承认其中一种是观看和视察的运作，或者正如人们可以称之为\Quote{凝视}，祂借此观看并俯瞰一切，洞察我们的思想，甚至凭借祂的默观能力进入那些不可见的事物，我们推测\Quote{神性}或\OriginalTerm{神性}{\GreekText{θε}{o}\GreekText{της}}的称呼来自\OriginalTerm{观看}{\GreekText{θ}{e}\GreekText{α}}或凝视，而那作为我们\OriginalTerm{观看者}{\GreekText{θεατ}{\GreekText{η}}\GreekText{ς}}或默观者的，按照习惯用法和圣经的教导，被称为\OriginalTerm{天主}{\GreekText{θε}{o}\GreekText{ς}}或天主。\newline{}如果有人承认观看和辨别是同一回事，并且那掌管万有的天主是宇宙的掌管者并被称为如此，让他考虑这个运作，判断它是否属于我们信仰的天主圣三中的一位，还是那能力扩展到三位。因为如果我们对\Quote{神性}或\OriginalTerm{神性}{\GreekText{θε}{o}\GreekText{της}}一词的解释是正确的，并且被看到的东西被称为\OriginalTerm{被观看之物}{\GreekText{θεατ}{a}}，而观看它们的那位被称为\OriginalTerm{天主}{\GreekText{θε}{o}\GreekText{ς}}或天主，那么圣三中的每一位都不能因该词所涉及的含义而被合理地排除在这一称呼之外。\newline{}因为圣经同样将观看的行为归于圣父、圣子和圣神。达味说：\Quote{天主，我们的护佑者，请观看}；由此我们得知观看是天主观念的一种适当运作，就天主被认知而言，因为他说：\Quote{天主，请观看}。\newline{}耶稣也看见那些谴责祂之人的思想，并质问凭祂自己的能力祂为何赦免人的罪？因为经上说：\Quote{耶稣看见他们的思想}。\newline{}关于圣神，伯多禄也对阿纳尼雅说：\Quote{为什么撒殚充满了你的心，叫你欺骗圣神？}\BibleRef{宗 5:3}，表明圣神是真实的见证者，知道阿纳尼雅暗中所行的事，并且通过祂，这秘密向伯多禄显明。因为阿纳尼雅偷窃自己的财产，暗自以为瞒过众人，隐藏了他的罪：但圣神同时就在伯多禄内，察觉了他的意图——这意图被拉向贪婪——并从祂自己赐给伯多禄看见秘密的能力，显然，如果祂不能看见隐藏的事，祂就不可能做到这一点。\par

但有人或许会说，我们论证的证据尚未触及这个问题。即使承认\Quote{\OriginalTerm{天主性}{Godhead}}是本性的共同称号，也并不能确定我们不应称\Quote{\OriginalTerm{众天主}{Gods}}；相反，通过这些论证，我们反倒被迫称\Quote{\OriginalTerm{众天主}{Gods}}：因为在人类的习惯中，我们发现不仅那些分享同一本性的人，甚至任何从事同一行业的人，当他们人数众多时，也不以单数称呼；正如我们说\Quote{许多演说家}、\Quote{测量员}、\Quote{农夫}或\Quote{制鞋匠}，其他所有情况也同样如此。的确，如果\OriginalTerm{天主性}{Godhead}是本性的称号，那么根据上述论证，更应当将三位包含在单数中，并因本性的不可分离性和不可分割性而称\Quote{独一天主}；但既然从前面的话中已经确定，\Quote{\OriginalTerm{天主性}{Godhead}}一词表示\OriginalTerm{行动}{operation}而非\OriginalTerm{本性}{nature}，那么从上述论证似乎得出相反的结论：我们因此更应当称那些在相同行动中被思考者为\Quote{\OriginalTerm{三位天主}{three Gods}}，如同他们说一个人会称\Quote{三位哲学家}或\Quote{三位演说家}，或任何其他由行业派生的名称，当参与同一行业的人不止一位时。\par

我在阐述此观点时费了一些力气，提出了代表反对者的论点，以便我们的决定因更详尽的反驳而得以加强，从而更加稳固。现在让我们继续我们的论证。\par

因我们在陈述中已在一定程度上表明，\Quote{\OriginalTerm{天主性}{Godhead}}一词并非表示本质，而是表示\OriginalTerm{运作}{operation}，或许有人可以合理地举出原因，为何就人而言，那些在同一追求中彼此分享的人，是以复数列举和提及，而另一方面，\OriginalTerm{神性}{Deity}却是以单数提及，称为一位天主和一位天主性，尽管三位\OriginalTerm{位格}{Persons}并未与\Quote{\OriginalTerm{天主性}{Godhead}}一词所表达的意义分离——我可以说，原因在于：人，即使多人从事同一形式的活动，也是各自单独工作于自己所承担的任务，其个人行动与从事同一职业的其他人没有参与。例如，假设有几位修辞学家，他们的追求是一致的，在多数情况下有相同的名称：但每一位从事这一追求的人都是独立工作，这一位为自己辩护，那一位也为自己辩护。因此，由于人在同一追求中的行动各有区别，他们被恰当地称为多数，因为每个人根据其运作的特殊性质，在自己的环境内与他人分离。但在神圣本性方面，我们并不类似地得知，圣父独自行任何圣子不合作的事，或者圣子有任何与圣神分离的特殊运作；相反，每一运作，凡由天主延伸至受造界，并按照我们对其多变的概念而命名的，都源于圣父，经由圣子而进行，并在圣神内完成。因此，源于运作的名称，并不根据实行者的数量而分割，因为每个受造物之行动并非分离和独特，而是凡发生的一切，无论是关于祂对我们的眷顾作为，还是关于宇宙的统治和构成，都是由三位的行动而发生，然而所发生的事并非三件事。我们可以通过一个实例来理解这一点。我说，从恩赐的源头本身，一切分享这恩宠的事物都获得了生命。那么，当我们询问这美好的恩赐是从何而来时，我们藉着圣经的指引发现，它来自圣父、圣子和圣神。然而，尽管我们提出三位位格和三个名字，我们并不认为我们被赋予了三个生命，分别来自每一位；而是同一生命由圣父在我们内运作，由圣子预备，并依靠圣神的意愿。既然天主圣三以类似我所说方式完成每一运作，并非按照位格的数量分别行动，而是有一善意的运动与安排，由圣父经由圣子传达给圣神（因为正如我们不将那些运作赋予同一生命者称为三个赋予生命者，同样也不将那些在同一良善中被默观者称为三个良善者，也不以其他属性按复数称呼他们）；同样，我们也不能将那位对自我和一切受造物施行这神圣和监管能力与运作，共同而不可分割地相互行动者，称为三位神。因为，当我们从圣经的话语中得知，宇宙的天主审判普世 \BibleRef{罗 3:6} 时，我们说祂经由圣子审判万有；又当我们听到，圣父不审判任何人时，我们并不认为圣经自相矛盾——（因为审判普世者经由祂的圣子而行，祂将一切审判交托给圣子；凡由\OriginalTerm{独生子}{Only-begotten}所作的一切，都归向圣父，所以祂自己既是审判万有者，又不审判任何人，因为祂已如我们所说，将一切审判交托给圣子，而圣子的一切审判都符合圣父的旨意；人们不能恰当地说他们是两位审判者，也不能说其中一位被排除在审判所隐含的权威和能力之外）；——同样，对于\Quote{\OriginalTerm{天主性}{Godhead}}一词，基督是天主的德能和天主的智慧，而那种监督和观看的德能，我们称为天主性，圣父经由独生子行使，而圣子藉着圣神完成一切德能，按照依撒意亚所说，以审判的神和焚烧的神施行审判 \BibleRef{依 4:4}，并且也藉着祂行动，符合福音中向犹太人所说的话。因为祂说：\Quote{我若靠着天主的神驱魔 \BibleRef{玛 12:28}}；祂以部分的描述包括了每一种行善的形式，因为行动的统一：源于运作的名称不能在许多之间分割，因为他们的共同运作结果是单一的。\par

既然监督和视察之权能的特性在圣父、圣子、圣神内是同一个，正如我们先前论证中所说的，由圣父如泉源发出，由圣子付诸实行，并藉圣神的能力成全其恩宠；而且既然任何行动都不按位格分离，不是由每一位单独完成，与我们默观中与之相连的那一位无关，反而对万物的所有眷顾、关心和监督，无论是对感官受造物还是超性自然的事物，以及保存现存之物、纠正错失之处、指导正确次序的能力，都是同一个，而非三个，确实是由天主圣三所指导，却未按照信德中所默观的位格数目以三重划分而割裂，以致每个行为，若单独默观，应只是圣父独自的工作，或圣子特有的工作，或圣神分别的工作；然而，正如宗徒所说：\BibleQuote{同一位圣神将祂的美好恩赐分别赐给每个人}\BibleRef{格前 12:11}，从圣神出发的善的运动并非没有起始——我们发现我们默观为先行于这运动的能力，即独生的天主，是万物的创造者；没有祂，任何存在物都不能获得其存在的开始：而同样，这同一善的源泉出自圣父的旨意。\par

因此，如果每一善事和每一善名，依赖那无始的能力和目的，在圣神的能力中，通过独生的天主，得以成全，没有时间或区分的标记（因为从圣父，经由圣子，到圣神的天主旨意运动中没有存在或设想的延迟）：而且如果 \OriginalTerm{神性}{Godhead} 也是善名和概念之一，那么将这名号分割为多个就是不恰当的，因为行动中存在的统一性阻止了复数列举。正如\BibleQuote{万民的救主，特别是信徒的救主}\BibleRef{弟前 4:10}，被宗徒称为一位，无人以此话论证圣子不拯救信徒，或救恩不经圣神介入就赐予领受者；但统御万有的天主是万民的救主，而圣子藉着圣神的恩宠成就救恩，然而圣经并未因此称他们为三位救主（虽然救恩被宣认出自天主圣三）：同样，他们也不按照赋予术语\Quote{神性}的意义被称为三位神，尽管前述称谓附着于天主圣三。\par

在我看来，为了当前论证的证明，没有必要绝对反对那些以我们不应将\Quote{神性}视为一种行动来反对我们的人。因为我们相信天主性是无限和不可理解的，我们对它没有任何理解，而是宣告这性质在各方面都被视为无限：而且绝对无限的不在某方面受限而在另一方面不受限，而是无限完全不受限制。因此，不受限制的当然也不受名称的限制。为了在天主性的情况下标志我们对无限的恒定理解，我们说神性超乎一切名号之上：而\Quote{神性}是一个名号。现在，同一事物不能同时是一个名号又被视为超乎一切名号之上。\par

但如果我们的对手喜欢说这术语的意义不是行动而是性质，我们将回到我们最初的论证，即习俗错误地应用性质的名号来表示众多：因为按照真实推理，任何性质在默观较大或较小的数目时，既不增加也不减少。因为只有那些在个体限制内被默观的事物才以相加的方式被列举。现在这种限制由身体外貌、大小、地点、形状和颜色的差异所标记，而撇开这些条件被默观的事物则不受由这些范畴形成的限制。不被如此限制的就不被列举，不被列举的就不能在众多中被默观。因为我们说金子，即使被切割成许多形状，仍是一个，并如此称呼，但我们说许多钱币或许多斯塔特，却没有发现金子的性质因斯塔特的数目而增多；因此，当金子以较大块状被默观时，无论是片状还是币状，我们称其为\Quote{多}，但我们不因其材料的众多而称之为\Quote{众多金子}——除非有人说有\Quote{许多金币}（例如大流克或斯塔特），在这种情况下，不是材料，而是钱币适用于数目意义：实际上，严格来说，我们不应称它们为\Quote{金子}，而应称为\Quote{金制的}。\par

因此，正如金色的斯塔特是多数的，但金子是一，同样，那些在人的本性中分别显示给我们的人，如\Person{伯多禄}、\Person{雅各伯}、\Person{若望}，是多数的，然而在他们内的人是一个。虽然圣经按复数意义扩展了用词，如说\Quote{人指着更大的起誓\BibleRef{希 6:16}}和\Quote{世人}，以及其他类似的短语，我们必须认识到，圣经是在使用流行言语的习惯，并未制定关于以某种方式使用词语的规则，也不是通过教导我们关于短语而说这些话，而是按照流行习惯使用词语，唯一的目的在于使话语对领受者有益，在表达方式上并不斤斤计较于那些不致因短语被理解的方式而产生害处的细节。\par

的确，从圣经中详细列出那些表达不精确的语句，以证明我所说的论断，将是一项冗长的任务；然而，当真理的任何部分有受损的危险时，我们在圣经的措辞中就再也找不到随意或漫不经心的用词了。为此，圣经允许以复数形式命名\Quote{人}，因为没有人会因这种修辞而在观念中误以为有多种人性，或认为表达这种本性的名称使用复数时，就表明有许多人性。但\Quote{天主}一词，圣经刻意只使用单数形式，以防通过\Quote{诸神}的复数含义而引入神圣本质中存在不同本性的观念。这就是为什么圣经说\BibleQuote{主我们的天主是唯一的主} \BibleRef{申 6:4}，并且以天主性的名称宣告独生天主，而没有将合一分裂为双重含义，以致称圣父和圣子为两位天主，尽管圣作者们宣告每一位都是天主。圣父是天主，圣子是天主，然而同一宣告中天主是唯一的，因为在天主性中，无论本性还是运作，都没有看到任何差异。因为（按照那些被误导者的想法）如果天主圣三的本性是不同的，那么数目会相应地扩展为多位天主，按照主体本质的差异而分裂。但既然神圣、单一、不变的本性为保持合一而拒绝本质上的任何差异，它自身就不接受多数的含义；正如它被称为单一的本性，它也用单数称呼其所有其他名称：\Quote{天主}、\Quote{善}、\Quote{圣}、\Quote{救主}、\Quote{义}、\Quote{审判者}，以及每一个可想象的神圣名称：至于这些名称是指本性还是运作，我们不予争论。\par

然而，若有人挑剔我们的论证，认为由于不承认本性的差异，导致位格的混合与混淆，我们将对此指控作如下答复：我们既承认本性的不变特征，也不否认就原因与被原因所产生的方面的区别——我们正是借此才领会一个位格与另一个位格的区别——即我们相信，一位是\Quote{原因}，另一位是\Quote{出于原因}的；并且，在出于原因的那一位中，我们又辨认出另一区别。因为一位是直接出于第一原因，另一位则是通过那直接出于第一原因的；如此，\Quote{独生子}的属性毫无疑问地存于圣子，而圣子的中介，在维护祂的独生子属性的同时，并不将圣神排除于祂与圣父在自然关系之外。\par

然而在谈论\Quote{原因}和\Quote{出于原因}时，我们并非用这些词语来指称本性（因为没有人会给出\Quote{原因}和\Quote{本性}的同一界说），而是表明存在方式的差异。因为当我们说一位是\Quote{有原因的}，另一位是\Quote{无原因的}时，我们并非用\Quote{原因}一词来划分本性，而只是表明圣子并非无受生而存在，圣父也并非由受生而存在：但我们首先必须相信某物存在，然后审视我们所信之对象的存在方式：因此，存在问题是一回事，存在方式问题是另一回事。说某物无受生而存在，是在表述其存在方式，但该短语并未指明存在的是什么。如果有人问一个农夫关于一棵树，它是被栽种的还是自己生长的，他回答说树没有被栽种或是栽种的结果，难道这个回答会表明树的本性吗？当然不会；当他说它如何存在时，他反而让本性的问题晦暗不明。同样，在另一种情况下，当我们得知祂是\Quote{无生的}，我们就学到了祂以何种方式存在，以及我们应如何恰当地想象祂存在，但祂\Emph{是什么}，我们并未从这个短语中听到。因此，当我们在天主圣三中承认这样的区分，即相信一位是\Quote{原因}，另一位\Quote{出于原因}时，我们就不能再被指责因本性的共同性而混淆了位格的定义。\par

这样，既然一方面原因的观念区分了天主圣三的位格，宣告一位无原因而存在，另一位出于原因；而另一方面，神圣本性在每个概念中都被理解为不变且不可分割的，因此我们有理由宣告天主性是一，天主是一，并用单数使用所有其他表达神圣属性的名称。\par

\SourceRecord{Translated by H.A. Wilson.；Nicene and Post-Nicene Fathers, Second Series；Vol. 5.；Edited by Philip Schaff and Henry Wace.；Buffalo, NY: Christian Literature Publishing Co.,；1893.；<http://www.newadvent.org/fathers/2905.htm>.}
